\documentclass{article}
\usepackage[utf8]{inputenc}
\usepackage{amsmath}
\usepackage{amsfonts}
\usepackage{amssymb}
\usepackage{graphicx}
\usepackage{hyperref}
\usepackage{geometry}
\usepackage{tcolorbox}
\geometry{a4paper, margin=1in}

\title{Module 01: Deep Dive into Linear Algebra for Machine Learning}
\author{Cybernauts 2026 Datathon Campaign}
\date{May 2026}

\begin{document}

\maketitle

\begin{abstract}
Linear Algebra is not just about solving for $x$. In Machine Learning, it is the fundamental framework for representing data, performing efficient computations, and understanding the geometric structure of high-dimensional spaces. This note provides an intuitive and comprehensive guide to the concepts you need to win the datathon.
\end{abstract}

\tableofcontents
\newpage

\section{The Hierarchy of Data Structures in ML}
In ML, we don't just see numbers; we see \textbf{tensors}. A tensor is simply a generalization of scalars, vectors, and matrices to any number of dimensions.

\begin{itemize}
    \item \textbf{Scalars (0D):} A single magnitude. 
    \textit{Example:} The learning rate $\alpha = 0.001$.
    \item \textbf{Vectors (1D):} An ordered array of numbers. 
    \textit{Intuition:} A single data point. If you are predicting house prices, a vector $\mathbf{x} = [3, 1500, 2]^T$ represents one house with 3 bedrooms, 1500 sqft, and 2 baths.
    \item \textbf{Matrices (2D):} A grid of numbers. 
    \textit{Intuition:} A spreadsheet or a dataset. Rows are samples; columns are features. 
    \item \textbf{Tensors (ND):} Higher-dimensional arrays. 
    \textit{Example:} A color image is a 3D tensor of shape $(Height, Width, Channels)$. A video is a 4D tensor adding the Time dimension.
\end{itemize}

\section{Fundamental Operations: Alignment and Similarity}

\subsection{Matrix Transpose ($\mathbf{X}^T$)}
Transposing a matrix flips it over its diagonal. Rows become columns.
\textit{Why it matters:} In ML, we often need to "align" dimensions. If you have a feature matrix $\mathbf{X}$ (Samples $\times$ Features) and you want to calculate the correlation between features, you often need $\mathbf{X}^T\mathbf{X}$.

\subsection{The Vector Dot Product ($\mathbf{u} \cdot \mathbf{v}$)}
The dot product is a single number: $\sum u_i v_i$.
\begin{tcolorbox}[title=The Geometric Intuition]
The dot product measures \textbf{Similarity}. If two vectors point in the same direction, their dot product is large. If they are perpendicular ($90^\circ$), it is zero. In ML, we use this to see how well a feature vector aligns with our model's weights.
\end{tcolorbox}
\textbf{Prediction Engine:} $\hat{y} = \mathbf{w}^T \mathbf{x} + b$. This is just a dot product plus a bias.

\subsection{Matrix Multiplication: Batch Processing}
Multiplying a matrix $\mathbf{A}$ by a matrix $\mathbf{B}$ is like performing multiple dot products at once. 
\textit{ML Perspective:} When we pass a whole batch of 32 images through a neural network layer, we aren't doing 32 separate calculations. We are doing one single \textbf{Matrix Multiplication}. This is why GPUs make ML so fast—they are "Matrix Multiplication Machines."

\section{Systems of Equations and the Normal Equation}

\subsection{Linear Systems: $\mathbf{Xw} = \mathbf{y}$}
Most of ML is trying to solve this equation. We have data $\mathbf{X}$ and targets $\mathbf{y}$, and we want to find the "importance" of each feature, $\mathbf{w}$. 
In most cases, we have more equations (samples) than unknowns (features). This is called an \textbf{overdetermined system}. There is usually no perfect solution that hits every point exactly, so we look for the "Best Fit."

\subsection{The Closed-Form Normal Equation}
To find the weights $\mathbf{w}$ that minimize the squared error, we use:
\begin{equation}
    \mathbf{w} = (\mathbf{X}^T \mathbf{X})^{-1} \mathbf{X}^T \mathbf{y}
\end{equation}
\textit{When to apply:} Use this for small-to-medium datasets where you want the \textbf{exact} best solution without using iterative Gradient Descent. It requires the matrix $(\mathbf{X}^T \mathbf{X})$ to be \textbf{Invertible}.

\section{The Matrix Inverse and Identity Matrix}
\begin{itemize}
    \item \textbf{Identity Matrix ($\mathbf{I}$):} The "number 1" of matrices. $\mathbf{AI} = \mathbf{A}$. It has 1s on the diagonal and 0s elsewhere.
    \item \textbf{Inverse ($\mathbf{A}^{-1}$):} The "reciprocal." $\mathbf{A} \mathbf{A}^{-1} = \mathbf{I}$. 
\end{itemize}
\textit{Warning:} Not all matrices have an inverse. If a matrix has redundant columns (perfectly correlated features), it is "singular" and cannot be inverted. This is why we remove highly correlated features in data cleaning!

\section{Eigen-everything: Variance and PCA}

\subsection{Eigenvalues and Eigenvectors}
An eigenvector $\mathbf{v}$ is a direction that stays constant during a transformation. The eigenvalue $\lambda$ is the factor by which it scales.
\begin{equation}
    \mathbf{A}\mathbf{v} = \lambda\mathbf{v}
\end{equation}

\subsection{The Covariance Matrix}
The covariance matrix $\mathbf{\Sigma}$ summarizes how features vary together. If we take the eigenvectors of this matrix, we find the \textbf{Principal Components}.

\subsection{Principal Component Analysis (PCA)}
PCA is the king of dimensionality reduction. 
\begin{enumerate}
    \item It finds the eigenvector with the largest eigenvalue. This is the direction of \textbf{Maximum Variance}.
    \item It projects the data onto these axes.
\end{enumerate}
\textit{Intuition:} If you have a 3D cloud of data that looks like a flat pancake, PCA will find the two directions that define the pancake and ignore the "thickness," reducing your 3D data to 2D without losing the main structure.

\section{Vector Norms for Regularization}
How do we stop a model from "cheating" by using massive weights to fit noise? We penalize the "size" of the weight vector $\mathbf{w}$.

\begin{itemize}
    \item \textbf{$L_1$ Norm (Manhattan):} $\|\mathbf{w}\|_1 = \sum |w_i|$. 
    \textit{Effect:} It forces some weights to become exactly \textbf{zero}. Great for \textbf{Feature Selection}.
    \item \textbf{$L_2$ Norm (Euclidean):} $\|\mathbf{w}\|_2 = \sqrt{\sum w_i^2}$. 
    \textit{Effect:} It keeps weights \textbf{small and balanced}. This is the standard "Weight Decay" used in almost all deep learning models to prevent overfitting.
\end{itemize}

\section{Practice Problems}
\textbf{P1:} You have a dataset $\mathbf{X}$ where Column 1 is "Height in cm" and Column 2 is "Height in inches." Why will the Normal Equation fail? \\
\textit{Answer: The columns are linearly dependent (perfectly correlated). The matrix $\mathbf{X}^T\mathbf{X}$ will be singular and non-invertible.}

\textbf{P2:} If you want to compress your data from 100 features to 10 features while keeping as much "information" as possible, which Linear Algebra tool do you use? \\
\textit{Answer: PCA, by selecting the 10 eigenvectors with the largest eigenvalues.}

\end{document}
